\chapter{Categories}

\begin{center}\itshape
A group can be described by its elements. It can also be described by the maps that enter or leave it. How much of the group does the second description remember?
\end{center}

Take a homomorphism from $\Z$ to a group $G$. Once you choose where $1$ goes, every other value is forced: $n$ goes to the $n$th power of that element. Thus the elements of $G$ can be recovered from the maps $\Z\to G$. The multiplication is visible too, if we allow maps from the free group on two generators. Maps carry more information than their notation first suggests.

Categories give us a language for asking which information is carried by maps. The examples are familiar: a free group, an abelianisation, the double dual of a vector space, and a function of two variables written as a family of functions of one variable. Each answers a different question. Their common feature is that the answer is specified by a rule about maps, rather than by a choice of coordinates.

Samuel Eilenberg and Saunders Mac Lane developed this language while working on group extensions and homology. In the early 1940s Mac Lane spoke at Michigan about groups $E$ containing an abelian group $A$, with quotient $E/A$ another prescribed abelian group $B$. Eilenberg recognised a connection with Steenrod's question about the homology of the complement of a $p$-adic solenoid. Their collaboration led to the 1945 paper ``General theory of natural equivalences''. Mac Lane's later account explains the order of invention: categories made it possible to define functors, and functors made it possible to define natural transformations.

Another conversation gave the subject one of its central results. Yoneda explained his ideas to Mac Lane in Paris; the discussion began at the Gare du Nord and continued on Yoneda's train. Yoneda's lemma says precisely how an object is remembered by its maps. Kan's 1958 work on adjoint functors then showed that many familiar constructions share a single existence-and-uniqueness pattern. The historical accounts are listed at the end of the chapter. We will meet the pattern first in sets, groups and vector spaces, where each arrow can be checked directly.

\section{Categories}

Suppose $f:A\to B$ and $g:B\to C$ are homomorphisms. Their composite is another homomorphism. We can compose three maps without deciding which pair to compose first, and the identity map changes nothing. A category retains exactly this information.

\begin{definition}\label{def:category}
A \emph{category} $\mathcal C$ has a class of \emph{objects}. For each pair $A,B$ it has a set $\operatorname{Hom}_{\mathcal C}(A,B)$, whose members are called \emph{morphisms} or \emph{arrows} and are written $f:A\to B$. For each triple of objects it has a composition operation
\[
\operatorname{Hom}(B,C)\times\operatorname{Hom}(A,B)\longrightarrow\operatorname{Hom}(A,C),\qquad(g,f)\longmapsto g\circ f.
\]
These data satisfy two requirements:
\begin{enumerate}[label=(\roman*)]
\item for composable arrows, $h\circ(g\circ f)=(h\circ g)\circ f$;
\item each object $A$ has an arrow $\mathrm{id}_A:A\to A$ such that $f\circ\mathrm{id}_A=f$ and $\mathrm{id}_B\circ f=f$ whenever $f:A\to B$.
\end{enumerate}
\end{definition}

The order in $g\circ f$ matters: apply $f$ first. Its target must be the source of $g$. You cannot compose two arrows merely because both occur in the same category.

There is no set of all sets, as Volume~7, Chapter~1 explains. This is why the objects may form a class while each individual collection $\operatorname{Hom}(A,B)$ is required to be a set. This requirement is called \emph{local smallness}. A category is \emph{small} if its objects and arrows together form sets. We allow large categories such as the category of sets. A natural-transformation family over a large category is initially a class-sized family; Yoneda's lemma will give a set parametrising the particular families used in its statement. No set of all sets is needed.

\begin{example}\label{ex:categories}
The following choices all satisfy the two composition rules.
\begin{enumerate}[label=(\alph*)]
\item In $\mathbf{Set}$ the objects are sets and the arrows are functions. In $\mathbf{Grp}$ they are groups and homomorphisms; in $\mathbf{Ab}$, abelian groups and homomorphisms. The categories $\mathbf{Ring}$ and $\mathbf{CRing}$ use rings with $1$ and maps preserving $1$, with commutativity required in the latter. The zero ring is allowed. In $R\text{-}\mathbf{Mod}$ the arrows are $R$-linear maps, and $\mathbf{Vect}_K$ is the case of vector spaces over a field $K$. In $\mathbf{Top}$ the arrows are continuous maps; in $\mathbf{Top}_*$ both spaces and maps have a chosen base point. In each case composition is ordinary composition of functions.
\item A partially ordered set $(P,\le)$ gives a category: there is one arrow $x\to y$ if $x\le y$, and no arrow otherwise. Reflexivity supplies identities. Transitivity supplies composition, and any two arrows with the same source and target are already equal.
\item A monoid $M$ gives a category with one object $\bullet$ and $\operatorname{Hom}(\bullet,\bullet)=M$. Composition is the monoid multiplication. For a group every arrow in this category has an inverse.
\item The objects of $\mathbf{Mat}_K$ are $0,1,2,\ldots$. An arrow $n\to m$ is an $m\times n$ matrix, and composition is matrix multiplication. The object $n$ records a dimension, rather than a particular vector space. Nevertheless its arrows act as linear maps $K^n\to K^m$. This example shows that objects need not themselves be sets equipped with extra operations.
\item Reverse every arrow of $\mathcal C$ to obtain its \emph{opposite category} $\mathcal C^{\mathrm{op}}$. Its objects are unchanged and
$\operatorname{Hom}_{\mathcal C^{\mathrm{op}}}(A,B)=\operatorname{Hom}_{\mathcal C}(B,A)$.
Composition is the old composition read in reverse order.
\end{enumerate}
\end{example}

An arrow $f:A\to B$ is an \emph{isomorphism} when there is $g:B\to A$ with $g\circ f=\mathrm{id}_A$ and $f\circ g=\mathrm{id}_B$. Its inverse is unique: if $g'$ also works, then
\[
g'=g'\circ\mathrm{id}_B=g'\circ f\circ g=\mathrm{id}_A\circ g=g.
\]
We write it as $f^{-1}$. Isomorphisms are bijections in $\mathbf{Set}$ and group isomorphisms in $\mathbf{Grp}$. In $\mathbf{Top}$ they are homeomorphisms. A continuous bijection need not have a continuous inverse, so forgetting which arrows are allowed can change the answer.

Cancellation gives two further notions. A \emph{monomorphism} $f:A\to B$ satisfies $f\circ g=f\circ h\Rightarrow g=h$ for every pair $g,h:X\to A$. An \emph{epimorphism} satisfies $g\circ f=h\circ f\Rightarrow g=h$ for every pair $g,h:B\to X$. In sets these are exactly injections and surjections. Exercise~\ref{ex:zq} shows why it is safer to keep the cancellation definitions in other categories: $\Z\to\Q$ is both a monomorphism and an epimorphism of rings, yet has no inverse.

\begin{think}
An arrow in a poset category can have an inverse only when both $x\le y$ and $y\le x$. What does antisymmetry then say? In $\mathbf{Mat}_K$, use the dimensions of the source and target to decide which arrows can be isomorphisms, including the object $0$.
\end{think}

\section{Functors}

Passing from a group to its underlying set also sends every group homomorphism to a function. The assignment respects composites and identities. A functor is an assignment with these two properties, whether or not it forgets anything.

\begin{definition}\label{def:functor}
A \emph{functor} $F:\mathcal C\to\mathcal D$ assigns an object $F(A)$ to every object $A$ and an arrow $F(f):F(A)\to F(B)$ to every arrow $f:A\to B$, subject to
\[
F(g\circ f)=F(g)\circ F(f),\qquad F(\mathrm{id}_A)=\mathrm{id}_{F(A)}.
\]
A \emph{contravariant functor} from $\mathcal C$ to $\mathcal D$ means a functor $\mathcal C^{\mathrm{op}}\to\mathcal D$. It sends $f:A\to B$ to $F(f):F(B)\to F(A)$, so
$F(g\circ f)=F(f)\circ F(g)$.
\end{definition}

The rule on arrows is part of the data. Giving only the object $F(A)$ does not yet give a functor.

\begin{example}\label{ex:functors}
Each construction below comes with its rule on arrows.
\begin{enumerate}[label=(\alph*)]
\item A \emph{forgetful functor} discards structure: $\mathbf{Grp}\to\mathbf{Set}$ forgets multiplication, $\mathbf{Ring}\to\mathbf{Ab}$ retains addition, and $\mathbf{Top}\to\mathbf{Set}$ forgets the topology. It leaves the underlying functions unchanged.
\item The \emph{abelianisation} of $G$ is $G^{\mathrm{ab}}=G/G'$, where $G'$ is the subgroup generated by commutators $aba^{-1}b^{-1}$ (Volume~3, Exercise~2.4). A homomorphism $f:G\to H$ sends each commutator to a commutator. It therefore induces $f^{\mathrm{ab}}:G^{\mathrm{ab}}\to H^{\mathrm{ab}}$, given by $gG'\mapsto f(g)H'$. The induced map of a composite is the composite of the induced maps.
\item For commutative rings, $R\mapsto R^\times$ and $R\mapsto\mathrm{GL}_n(R)$ define functors to $\mathbf{Grp}$. Apply a ring map to a unit or to every matrix entry. Applying it to $MM^{-1}=I$ proves that the image matrix remains invertible.
\item The constructions $\pi_1:\mathbf{Top}_*\to\mathbf{Grp}$ and $H_n:\mathbf{Top}\to\mathbf{Ab}$ from Volume~9, Chapters~5 and~6 are functors. The identities for induced maps proved there, such as $H_n(g\circ f)=H_n(g)\circ H_n(f)$, are exactly the functor rules.
\item The dual $V^*=\operatorname{Hom}(V,K)$ gives a contravariant functor. A linear map $f:V\to W$ sends a form $\psi\in W^*$ to $f^*\psi=\psi\circ f\in V^*$. Indeed $(g\circ f)^*=f^*\circ g^*$.
\item Fix an object $A$. The functor $h^A=\operatorname{Hom}(A,-)$ sends an object $B$ to the set of arrows $A\to B$, and an arrow $g:B\to C$ to the function $f\mapsto g\circ f$. The contravariant functor $h_A=\operatorname{Hom}(-,A)$ sends $g:B\to C$ to $f\mapsto f\circ g$, from $\operatorname{Hom}(C,A)$ to $\operatorname{Hom}(B,A)$.
\item Between poset categories, a functor is an order-preserving map. Between groups regarded as one-object categories, it is a group homomorphism: the composition rule says that it preserves multiplication.
\end{enumerate}
\end{example}

A functor may lose information, but it cannot destroy the equations defining an inverse.

\begin{proposition}\label{prop:functoriso}
If $f$ is an isomorphism, then $F(f)$ is an isomorphism and its inverse is $F(f^{-1})$.
\end{proposition}

\begin{proof}
Apply $F$ to $f^{-1}\circ f=\mathrm{id}_A$. The functor rules give
\[
F(f^{-1})\circ F(f)=\mathrm{id}_{F(A)}.
\]
Applying $F$ to $f\circ f^{-1}=\mathrm{id}_B$ gives the inverse identity in the other order. Both identities are needed, and together they prove the assertion.
\end{proof}

This lets us disprove isomorphism in one category by computing in another. If a functor gives non-isomorphic groups from two spaces, the spaces cannot have been homeomorphic. It also lets us rule out maps satisfying particular equations.

\begin{example}[Brouwer, functorially]\label{ex:brouwer}
Let $i:S^1\to D^2$ include the boundary circle into the closed disc. A \emph{retraction} would be a continuous map $r:D^2\to S^1$ satisfying $r\circ i=\mathrm{id}_{S^1}$. Applying $H_1$ would give
\[
H_1(r)\circ H_1(i)=\mathrm{id}_{H_1(S^1)}.
\]
But Volume~9, Chapter~6 computes $H_1(D^2)=0$ and $H_1(S^1)\cong\Z$. Thus the left side factors through the zero group and is zero, whereas the right side is not zero. No such retraction exists. As in Volume~9, a continuous self-map of $D^2$ without a fixed point would give a retraction by extending the ray from the image through the original point to the boundary. Hence every continuous self-map of the disc has a fixed point.
\end{example}

\section{Natural transformations}

A functor compares objects. We also need to compare two functors. For example, each group has a quotient map to its abelianisation. The maps for different groups are compatible with every homomorphism between those groups. This compatibility is what the word \emph{natural} requires.

\begin{definition}\label{def:nat}
For functors $F,G:\mathcal C\to\mathcal D$, a \emph{natural transformation} $\eta:F\Rightarrow G$ chooses an arrow $\eta_A:F(A)\to G(A)$ for every object $A$, with
\[
G(f)\circ\eta_A=\eta_B\circ F(f)\qquad(f:A\to B).
\]
Thus the square in Figure~\ref{fig:universal} commutes for every arrow $f$. The transformation is a \emph{natural isomorphism} if every $\eta_A$ is an isomorphism. In that case the inverses form a natural transformation $G\Rightarrow F$.
\end{definition}

For the last assertion, start from the displayed equality and compose on the left with $\eta_B^{-1}$ and on the right with $\eta_A^{-1}$. The result is the naturality equation for the inverses. Naturality does not mean that an arbitrary choice has happened to work for each object separately. It says that all choices commute with all allowed arrows.

\begin{example}\label{ex:nattrans}
Two examples show how to check the square on elements.
\begin{enumerate}[label=(\alph*)]
\item The quotient maps $\pi_G:G\to G^{\mathrm{ab}}$ are natural. Regard abelianisation as a functor to $\mathbf{Grp}$ by including $\mathbf{Ab}$. For $f:G\to H$ and $g\in G$,
\[
f^{\mathrm{ab}}(\pi_G(g))=f(g)H'=\pi_H(f(g)).
\]
Both routes send an element to the same coset.
\item The determinants $\det_R:\mathrm{GL}_n(R)\to R^\times$ are natural in the commutative ring $R$. A ring homomorphism $\varphi:R\to S$ satisfies $\det(\varphi(M))=\varphi(\det M)$: the determinant is a sum of signed products of entries, with integer coefficients, and $\varphi$ respects every operation in that sum.
\end{enumerate}
\end{example}

\begin{think}
Fix an arrow $f:H\to K$. For each $X$, precomposition gives $-\circ f:\operatorname{Hom}(K,X)\to\operatorname{Hom}(H,X)$. Check the naturality equation for an arrow $X\to Y$. Why does the resulting transformation $h^K\Rightarrow h^H$ run opposite to $f$?
\end{think}

The double dual provides a useful test of this language. Volume~8, Theorem~A.8 proves that evaluation identifies a finite-dimensional space with its double dual. Evaluating a form at a vector does not ask for a basis. A direct identification with the single dual does.

\begin{theorem}[the double dual is natural]\label{thm:bidualnat}
On the category $\mathbf{FinVect}_K$ of finite-dimensional vector spaces, evaluation
\[
\iota_V:V\longrightarrow V^{**},\qquad\iota_V(v)(\varphi)=\varphi(v),
\]
is a natural isomorphism from the identity functor to the functor $V\mapsto V^{**}$, $f\mapsto f^{**}$.
\end{theorem}

\begin{proof}
The cited double-dual theorem makes each $\iota_V$ an isomorphism, including when $V=0$. It remains to check compatibility with $f:V\to W$. Two elements of $W^{**}$ are equal if they have the same value on every $\psi\in W^*$. For $v\in V$,
\[
\begin{aligned}
f^{**}(\iota_V(v))(\psi)
 &=\iota_V(v)(f^*\psi)\\
 &=(\psi\circ f)(v)\\
 &=\psi(f(v))=\iota_W(f(v))(\psi).
\end{aligned}
\]
Thus $f^{**}\circ\iota_V=\iota_W\circ f$, the required naturality equation.
\end{proof}

Evaluation still defines a natural map $V\to V^{**}$ in infinite dimension, but it is no longer onto; Volume~8, Appendix~A supplies the obstruction. Finite dimension is needed for the word \emph{isomorphism}, not for the commuting square.

The single dual reverses arrows, so first we must put both assignments in the same direction. Let $\mathbf{Iso}_K$ have finite-dimensional spaces as objects and only linear isomorphisms as arrows. For such an arrow $f:V\to W$, use $(f^*)^{-1}:V^*\to W^*$. These maps make the dual a covariant functor on $\mathbf{Iso}_K$. A natural family $\eta_V:V\to V^*$ would satisfy
\begin{equation}\label{eq:dualnat}
\eta_V=f^*\circ\eta_W\circ f.
\end{equation}
Here we multiplied the usual naturality equation by $f^*$ to put it in a convenient form.

\begin{theorem}[$V$ is not naturally its dual]\label{thm:notnatural}
Suppose $K$ has more than three elements. There is no family of isomorphisms $\eta_V:V\to V^*$ on $\mathbf{Iso}_K$ satisfying \eqref{eq:dualnat}. For each nonzero $V$, scalar isomorphisms $f=\lambda\,\mathrm{id}_V$ already prevent such a choice.
\end{theorem}

\begin{proof}
Put $W=V$ and $f=\lambda\,\mathrm{id}_V$ for $\lambda\in K^\times$. Since $f^*=\lambda\,\mathrm{id}_{V^*}$, equation~\eqref{eq:dualnat} becomes
\[
\eta_V=\lambda^2\eta_V.
\]
For nonzero $V$, an isomorphism $\eta_V$ is a nonzero linear map. The equality therefore forces $\lambda^2=1$. If naturality held for all scalar isomorphisms, every nonzero element of $K$ would be a root of $X^2-1$. A degree-two polynomial over a field has at most two roots, giving $|K^\times|\le2$, contrary to $|K|>3$.
\end{proof}

The zero-dimensional space causes no contradiction: its unique map to its dual is an isomorphism. The field hypothesis also matters. A plane over $\mathbb F_2$ has a natural self-duality, as Exercise~\ref{ex:f2plane} shows. Over larger fields a bilinear or inner-product structure can give an identification with a dual, but that identification respects only the maps preserving the extra structure. Naturality always depends on the arrows we permit.

\section{Universal properties}

A product set contains pairs. More useful than this description is the fact that a map into the product is exactly a pair of maps into its factors. This rule continues to make sense when objects have no elements that we can inspect.

\begin{definition}\label{def:initial}
An object $I$ is \emph{initial} if for every object $X$ there is exactly one arrow $I\to X$. An object $T$ is \emph{terminal} if for every $X$ there is exactly one arrow $X\to T$.
\end{definition}

The empty set is initial in $\mathbf{Set}$: there is one empty function to every set. A one-point set is terminal. The trivial group is both initial and terminal in $\mathbf{Grp}$. In $\mathbf{Ring}$, $\Z$ is initial because a map must send $n$ to $n\cdot1$; the zero ring is terminal. For a poset, an initial object is a least element and a terminal object a greatest element. Existence is part of the question: a poset need not have either.

\begin{proposition}\label{prop:initialunique}
Any two initial objects are joined by exactly one isomorphism. The corresponding assertion holds for terminal objects.
\end{proposition}

\begin{proof}
Let $I,I'$ be initial. Their properties give one arrow $u:I\to I'$ and one arrow $v:I'\to I$. Initiality of $I$ says there is only one endomorphism $I\to I$, so $v\circ u=\mathrm{id}_I$. The same argument at $I'$ gives $u\circ v=\mathrm{id}_{I'}$. Hence $u$ is an isomorphism, and it is unique because there was only one arrow $I\to I'$. Reversing arrows proves the terminal assertion.
\end{proof}

This is the source of the phrase ``unique up to a unique isomorphism''. When an object carries extra maps, the isomorphism must respect those maps. A product makes this qualification visible.

\begin{definition}\label{def:product}
A \emph{product} of $A,B$ is an object $P$ with arrows $p_A:P\to A$ and $p_B:P\to B$ such that, for every pair $f:X\to A$, $g:X\to B$, there is exactly one $h:X\to P$ with
\[
p_A\circ h=f,\qquad p_B\circ h=g.
\]
We write $P=A\times B$ and $h=(f,g)$. A \emph{coproduct} is the arrow-reversed construction: an object $A\sqcup B$ with maps $i_A:A\to A\sqcup B$ and $i_B:B\to A\sqcup B$ such that every $f:A\to X$, $g:B\to X$ factors uniquely as $f=h\circ i_A$, $g=h\circ i_B$.
\end{definition}

To see the uniqueness of products, make a category of triples $(X,f,g)$. An arrow from $(X,f,g)$ to $(X',f',g')$ is $k:X\to X'$ with $f'k=f$ and $g'k=g$. A product is terminal in this category. Proposition~\ref{prop:initialunique} therefore gives a unique isomorphism between two products \emph{commuting with both projections}. It does not say that the underlying object has only one automorphism.

\begin{figure}[ht]
\centering
\fig{ch01-universal}
\caption{Left: the universal property of the product. Every pair of arrows $f$, $g$ out of $X$ factors through one arrow $h=(f,g)$ into $A\times B$. Right: the naturality square of a
natural transformation $\eta\colon F\Rightarrow G$. Going right then down equals going down then right.}
\label{fig:universal}
\end{figure}

\begin{example}\label{ex:products}
The same definition produces several different constructions.
\begin{enumerate}[label=(\alph*)]
\item For sets, groups, rings and modules, take the cartesian product and use operations component by component. A map $h$ is forced to be $h(x)=(f(x),g(x))$, and this formula preserves the relevant operations.
\item For topological spaces, put the product topology on $A\times B$. A map to this space is continuous precisely when both components are continuous. Indeed inverse images of basic open sets $U\times V$ are $f^{-1}(U)\cap g^{-1}(V)$. The projections are continuous, and this is the coarsest topology with that property.
\item For a poset, the product is the greatest lower bound $\inf(x,y)$ and the coproduct is the least upper bound, when they exist. In positive integers ordered by divisibility these are $\gcd(x,y)$ and $\operatorname{lcm}(x,y)$.
\item For sets the coproduct is disjoint union, so equal-looking elements in different summands stay distinct. For modules it is $A\oplus B$: the injections are $a\mapsto(a,0)$ and $b\mapsto(0,b)$, and the extension is $h(a,b)=f(a)+g(b)$. For two modules the product and coproduct thus have the same underlying module, with different specified maps. For groups, $f(a)g(b)$ need not be a homomorphism from $A\times B$, because the two images need not commute. The coproduct is instead the free product $A*B$ of Volume~9, Proposition~5.17.
\end{enumerate}
\end{example}

Figure~\ref{fig:v10new-ch01-divisors} is a small case in which the arrows can all be read at once. For $4$ and $6$, a common incoming object is a common divisor, so the universal incoming object is $2$. A common outgoing object is a common multiple, so its universal choice is $12$.

\begin{figure}[ht]
\centering
\fig{v10-new-ch01-divisors}
\caption{The divisors of $12$, ordered by divisibility. Each upward edge is an arrow, and longer arrows are obtained by composition. The greatest common source of $4$ and $6$ is $2$, their product in this category. Their least common target is $12$, their coproduct.}
\label{fig:v10new-ch01-divisors}
\end{figure}

Maps out of an abelianisation have a similarly precise rule.

\begin{proposition}[universal property of abelianisation]\label{prop:abuniv}
For a group $G$ and an abelian group $A$, every homomorphism $f:G\to A$ has a unique factorisation $f=\bar f\circ\pi_G$, where $\bar f:G^{\mathrm{ab}}\to A$ is a homomorphism and $\pi_G:G\to G/G'$ is the quotient map.
\end{proposition}

\begin{proof}
The image of each commutator is
\[
f(aba^{-1}b^{-1})=f(a)f(b)f(a)^{-1}f(b)^{-1}=1,
\]
because $A$ is abelian. Thus $G'\subseteq\ker f$. The quotient factorisation theorem of Volume~3, Chapter~2 gives a well-defined map $\bar f(gG')=f(g)$, and it is a homomorphism. Its composite with $\pi_G$ is $f$. Every coset has a representative $g$, so this formula also proves uniqueness.
\end{proof}

In the category of homomorphisms from $G$ to abelian groups, $(G^{\mathrm{ab}},\pi_G)$ is initial. Any other proposed answer must receive exactly one compatible map from it. This describes what ``best abelian approximation'' means without a vague appeal to closeness.

The free module has the same kind of description. For a set $S$, let $R^{(S)}$ consist of finitely supported functions $S\to R$, with basis vectors $e_s$. Any function $x:S\to M$ into a module extends to exactly one linear map,
\[
\sum_{s\in S}c_se_s\longmapsto\sum_{s\in S}c_sx(s).
\]
Both sums are finite, even when $S$ is infinite. The formula proves existence; linearity forces the formula and so proves uniqueness. For $S=\varnothing$ the free module is $0$, agreeing with its initial property.

\section{The Yoneda lemma}

Fix a functor $F:\mathcal C\to\mathbf{Set}$ and an object $A$. If you choose $x\in F(A)$, every arrow $f:A\to B$ transports it to $F(f)(x)\in F(B)$. These transport rules are compatible with composition. Yoneda's lemma says that every natural way to turn an arrow out of $A$ into an element of $F(B)$ comes from one such $x$.

\begin{theorem}[Yoneda's lemma]\label{thm:yoneda}
For an object $A$ and a functor $F:\mathcal C\to\mathbf{Set}$, evaluation at the identity gives a bijection
\[
\Phi:\operatorname{Nat}(h^A,F)\longrightarrow F(A),\qquad
\Phi(\alpha)=\alpha_A(\mathrm{id}_A).
\]
The transformation corresponding to $x\in F(A)$ has components
$\alpha^x_B(f)=F(f)(x)$ for every arrow $f:A\to B$.
\end{theorem}

\begin{proof}
\emph{Build a transformation from an element.} Given $x\in F(A)$, define $\alpha^x_B$ by the stated formula. For $g:B\to C$ and $f:A\to B$,
\[
\alpha^x_C(g\circ f)=F(g\circ f)(x)=F(g)(F(f)(x))=F(g)(\alpha^x_B(f)).
\]
This is the naturality equation, so the components really form a transformation.

\emph{Recover the chosen element.} At $B=A$ and $f=\mathrm{id}_A$, the formula gives $\alpha^x_A(\mathrm{id}_A)=F(\mathrm{id}_A)(x)=x$. Hence $\Phi(\alpha^x)=x$.

\emph{Recover any transformation.} Now start from $\alpha$ and set $x=\alpha_A(\mathrm{id}_A)$. Naturality for the arrow $f:A\to B$, evaluated at $\mathrm{id}_A$, yields
\[
\alpha_B(f)=\alpha_B(h^A(f)(\mathrm{id}_A))
 =F(f)(\alpha_A(\mathrm{id}_A))=F(f)(x).
\]
Thus every component of $\alpha$ agrees with $\alpha^x$. The two constructions are inverse, proving the bijection.
\end{proof}

The identity is the useful test element: every $f:A\to B$ is obtained from it by postcomposition. Figure~\ref{fig:v10new-ch01-yoneda} shows the exact square used in the last step. The lemma is not asserting that arbitrary families of maps are determined at one object. Naturality is what forces the other components.

\begin{figure}[ht]
\centering
\fig{v10-new-ch01-yoneda}
\caption{Yoneda's determination of a natural transformation. The upper arrow sends $\mathrm{id}_A$ to $f$; the lower arrow sends $x=\alpha_A(\mathrm{id}_A)$ to $F(f)(x)$. Commutativity forces $\alpha_B(f)=F(f)(x)$ for every $f$.}
\label{fig:v10new-ch01-yoneda}
\end{figure}

Reading the same proof in $\mathcal C^{\mathrm{op}}$ gives the contravariant form: for contravariant $F$, transformations $h_A\Rightarrow F$ are parametrised by $F(A)$. This reversal explains the two directions in the following consequence.

\begin{corollary}[an object is known by its maps]\label{cor:yonedaemb}
Every natural transformation $h^A\Rightarrow h^B$ is precomposition $-\circ u$ by exactly one arrow $u:B\to A$. Consequently $h^A\cong h^B$ naturally if and only if $A\cong B$. In the contravariant version, transformations $h_A\Rightarrow h_B$ are exactly postcomposition $u\circ-$ by arrows $u:A\to B$, and $h_A\cong h_B$ naturally if and only if $A\cong B$.
\end{corollary}

\begin{proof}
Set $F=h^B$ in Yoneda's lemma. The parameter set is $h^B(A)=\operatorname{Hom}(B,A)$, and the transformation of $u$ sends $f:A\to X$ to $f\circ u:B\to X$. Follow this with the transformation belonging to $v:C\to B$: the result sends $f$ to $f\circ u\circ v$. Thus composition of transformations corresponds to composition of the representing arrows, with their direction reversed.

If the transformation of $u:B\to A$ has inverse represented by $u':A\to B$, their two composites correspond to $u\circ u'=\mathrm{id}_A$ and $u'\circ u=\mathrm{id}_B$. Hence $u$ is an isomorphism. Conversely an isomorphism $u$ supplies inverse transformations by precomposition with $u$ and $u^{-1}$. Apply this argument in the opposite category for the last assertion.
\end{proof}

A functor naturally isomorphic to $h^A$ is called \emph{representable}, with \emph{representing object} $A$. The corollary says that a representing object is determined up to isomorphism. This gives a practical strategy: find an object whose maps encode the data in a functor, then turn natural transformations into individual elements.

\begin{example}[Cayley's theorem]\label{ex:cayley}
Regard a group $G$ as a one-object category. A functor $G\to\mathbf{Set}$ gives a set $X$ and permutations $F(g)$ with $F(gh)=F(g)F(h)$ and $F(1)=\mathrm{id}$; this is a $G$-action. A natural transformation is precisely an equivariant map.

The representable $h^\bullet$ is the left action of $G$ on itself. Yoneda identifies its equivariant self-maps with $G$: a map commuting with all left multiplications must send $x$ to $xu$, where $u$ is its value at $1$. Indeed $\alpha(x)=\alpha(x\cdot1)=x\alpha(1)$. Meanwhile left multiplication itself gives a homomorphism $G\to\operatorname{Sym}(G)$. It is injective because the permutation $x\mapsto gx$ sends $1$ to $g$, so different $g$ give different permutations. This is Cayley's theorem, now seen through the representable action.
\end{example}

\begin{example}[the natural operations on groups]\label{ex:natops}
Suppose an operation $w_G:G\to G$ is assigned to every group and commutes with every homomorphism: $f(w_G(g))=w_H(f(g))$. The components need only be functions, not group homomorphisms. Such families are the transformations $U\Rightarrow U$ for the forgetful functor $U:\mathbf{Grp}\to\mathbf{Set}$.

Evaluation at $1$ identifies $\operatorname{Hom}(\Z,G)$ naturally with $U(G)$, so $U\cong h^{\Z}$. Yoneda therefore gives
\[
\operatorname{Nat}(U,U)\cong\operatorname{Nat}(h^{\Z},U)\cong U(\Z)=\Z.
\]
For the integer $n$, the corresponding operation sends $g=f(1)$ to $f(n)=g^n$. These powers, including the constant identity operation for $n=0$ and inversion for $n=-1$, exhaust the possibilities. A proposed operation defined for every group can be tested on one generator of $\Z$. Exercise~\ref{ex:words} makes the same test with two generators.
\end{example}

\begin{think}
A unital ring map $\Z[X]\to R$ is determined by the image of $X$, which may be any element of $R$. Thus the underlying-set functor on commutative rings is represented by $\Z[X]$. Apply Yoneda: which operations on one ring element commute with every unital homomorphism? Compare $r\mapsto r^2+3r$ with an operation involving division by $2$.
\end{think}

\section{Adjunctions}

To specify a linear map from a free module, we specify the images of its basis elements. To specify a map from an abelianisation to an abelian group, we specify a homomorphism from the original group. These correspondences change the source and target together. An adjunction records their common pattern.

\begin{definition}\label{def:adjunction}
Functors $F:\mathcal C\to\mathcal D$ and $G:\mathcal D\to\mathcal C$ are \emph{adjoint}, written $F\dashv G$, if there are bijections
\[
\varphi_{A,B}:\operatorname{Hom}_{\mathcal D}(F(A),B)
 \longrightarrow\operatorname{Hom}_{\mathcal C}(A,G(B))
\]
natural in both variables. Explicitly, for $a:A'\to A$, $b:B\to B'$ and $h:F(A)\to B$,
\[
\varphi_{A',B'}(b\circ h\circ F(a))
 =G(b)\circ\varphi_{A,B}(h)\circ a.
\]
We call $F$ the \emph{left adjoint} and $G$ the \emph{right adjoint}.
\end{definition}

The naturality condition prevents us from choosing unrelated bijections for different pairs $(A,B)$. Precomposing or postcomposing a map must give the same answer before or after the correspondence.

\begin{example}\label{ex:adjunctions}
These familiar constructions provide adjunctions.
\begin{enumerate}[label=(\alph*)]
\item The free-module functor $S\mapsto R^{(S)}$ is left adjoint to the forgetful functor $U:R\text{-}\mathbf{Mod}\to\mathbf{Set}$. The bijection restricts a linear map to the basis $S$. Conversely, a function on $S$ extends by finite linear combinations. The free-group functor of Volume~9, Chapter~5 and the free commutative ring $S\mapsto\Z[X_s:s\in S]$ obey the same rule.
\item Abelianisation is left adjoint to the inclusion $\mathbf{Ab}\to\mathbf{Grp}$. Proposition~\ref{prop:abuniv} supplies the bijection $\bar f\mapsto\bar f\circ\pi_G$ between homomorphisms $G^{\mathrm{ab}}\to A$ and $G\to A$.
\item Fix a set $B$ and write $C^B$ for the set of functions $B\to C$. A function $h:A\times B\to C$ corresponds to $a\mapsto(b\mapsto h(a,b))$, a function $A\to C^B$. Evaluation, $(a,b)\mapsto k(a)(b)$, reverses this correspondence. Therefore $-\times B\dashv(-)^B$. This rearrangement of a two-variable function is called \emph{currying}, after Haskell Curry. It also works for empty sets: the unique empty function accounts for the boundary cases.
\item Giving a set the discrete topology is left adjoint to forgetting a topology, since every function from a discrete space is continuous. Giving it the indiscrete topology is right adjoint to forgetting, since every function into an indiscrete space is continuous.
\item For posets an adjunction is the equivalence $F(x)\le y\iff x\le G(y)$. For a finite Galois extension $L/K$ with group $\Gamma$, take the subgroup poset and the opposite of the intermediate-field poset. The maps $H\mapsto L^H$ and $E\mapsto\operatorname{Gal}(L/E)$ satisfy
\[
E\subseteq L^H\quad\Longleftrightarrow\quad H\subseteq\operatorname{Gal}(L/E).
\]
Both assertions say that every element of $H$ fixes every element of $E$. This is the Galois connection behind the terminology. The fundamental theorem of Galois theory (Volume~3, Chapter~10) says that in this finite Galois case the correspondence is an equivalence of these posets.
\end{enumerate}
\end{example}

The maps allowed by an adjunction determine either side once the other is known. Yoneda makes this claim precise.

\begin{theorem}[uniqueness of adjoints]\label{thm:adjunique}
If $G$ and $G'$ are right adjoints to the same functor $F$, then there is a natural isomorphism $G\Rightarrow G'$.
\end{theorem}

\begin{proof}
For each $B$, combine the two adjunction bijections to obtain
\[
\theta_{A,B}:\operatorname{Hom}(A,G(B))
 \longrightarrow\operatorname{Hom}(F(A),B)
 \longrightarrow\operatorname{Hom}(A,G'(B)).
\]
This is a natural isomorphism of contravariant functors in $A$. By Corollary~\ref{cor:yonedaemb}, it has the form $t\mapsto u_B\circ t$ for an isomorphism $u_B:G(B)\to G'(B)$. In particular $u_B$ is the image of $\mathrm{id}_{G(B)}$.

We must still check that the $u_B$ fit together. For $b:B\to B'$, naturality of $\theta$ in $B$ gives
\[
\theta_{G(B),B'}(G(b))=G'(b)\circ u_B.
\]
Its description by postcomposition with $u_{B'}$ gives the same left side as $u_{B'}\circ G(b)$. Hence $G'(b)\circ u_B=u_{B'}\circ G(b)$ for every $b$. This is the required naturality, and all components are isomorphisms.
\end{proof}

An adjunction also transports a universal property. A right adjoint, for instance, must take a product to a product with the induced projections.

\begin{theorem}[right adjoints preserve products]\label{thm:rapl}
Suppose $F\dashv G$ and $B_1\times B_2$ is a product in $\mathcal D$, with projections $p_1,p_2$. Then $G(B_1\times B_2)$, equipped with $G(p_1),G(p_2)$, is a product of $G(B_1),G(B_2)$ in $\mathcal C$.
\end{theorem}

\begin{proof}
Let $\psi=\varphi^{-1}$ be the inverse adjunction bijection. Given $f_j:A\to G(B_j)$ for $j=1,2$, transpose them to $\psi(f_j):F(A)\to B_j$. The product in $\mathcal D$ gives a unique $k:F(A)\to B_1\times B_2$ with $p_jk=\psi(f_j)$. Set $h=\varphi(k)$. Naturality in the target gives
\[
G(p_j)\circ h=\varphi(p_j\circ k)=f_j.
\]
Thus $h$ has the required projections.

For uniqueness, suppose $h'$ has the same projections. Naturality gives $\varphi(p_j\circ\psi(h'))=G(p_j)\circ h'=f_j$, so $p_j\circ\psi(h')=\psi(f_j)$. Uniqueness for the original product forces $\psi(h')=k$, and applying $\varphi$ yields $h'=h$.
\end{proof}

The forgetful functors from groups, rings, modules and spaces to sets have the left adjoints just described. They therefore preserve products: the underlying set of the product is a cartesian product. Reversing arrows shows that left adjoints preserve coproducts (Exercise~\ref{ex:leftadj}). In particular,
\[
(A*B)^{\mathrm{ab}}\cong A^{\mathrm{ab}}\oplus B^{\mathrm{ab}}.
\]
We have derived an identity about noncommuting words from a rule about maps, without rearranging a word.

\begin{beauty}
Mac Lane wrote that ``adjoint functors arise everywhere''. The examples in this book give that sentence a concrete meaning. Tensor products in Chapter~3 are left adjoints to $\operatorname{Hom}$; localisation in Chapter~5 also has a universal map; and induction and restriction for finite-group representations are adjoint in Chapter~9. The character identity called Frobenius reciprocity is one expression of the last correspondence.

There is also a connection with programming, recorded in the work of Moggi and Wadler listed below. An adjunction gives a composite functor $G\circ F$ with maps that insert and combine two applications of it; this additional structure is called a \emph{monad}. For the free-monoid adjunction, $G\circ F$ forms finite lists, insertion makes a one-entry list, and combination flattens a list of lists. We will not use monads in any proof. The point is that the same rule about maps can organise algebraic constructions and operations on data.
\end{beauty}

\section{What we proved, and what comes next}

The homomorphisms $\Z\to G$ recover the elements of $G$. Yoneda explains why this is more than an isolated trick: a natural transformation out of a representable functor is fixed by its value on the identity. Compatible maps into, or out of, an object remember that object up to isomorphism. Without compatibility, merely counting the maps would not say this.

Universal properties identify the particular maps that matter. For a product we specify projections; for an abelianisation we specify the quotient map. Their uniqueness is uniqueness with those maps attached. Adjunctions let the source and target vary together, and their natural bijections transfer constructions such as products. When using a universal property, ask both questions: does the desired arrow exist, and what forces it to be unique?

The next chapter puts coordinates back on the table. A bilinear operation on two vector spaces can be recorded as a three-way array. For matrix multiplication, decomposing that array into simple pieces counts the multiplications in an algorithm. Chapter~3 will then remove the coordinates again and define the tensor product of modules by its maps.

\section{Exercises}

\begin{exercise}[$\star$]\label{ex:zq}
Use the cancellation definitions to prove that the inclusion $\Z\to\Q$ is both a monomorphism and an epimorphism in $\mathbf{Ring}$, although it is not an isomorphism. Decide whether the same function is an epimorphism in $\mathbf{Set}$, and explain the difference.
\end{exercise}

\begin{exercise}[$\star$]\label{ex:nofieldproduct}
Work in the category of fields and unital ring homomorphisms. Prove that $\mathbb F_2$ and $\mathbb F_3$ admit neither a product nor a coproduct. Identify which of the required structural maps cannot exist.
\end{exercise}

\begin{exercise}[$\star\star$]\label{ex:words}
Let $U:\mathbf{Grp}\to\mathbf{Set}$ forget the group structure, and let $U^2(G)=U(G)\times U(G)$, with arrows acting in both coordinates.
\begin{enumerate}[label=(\alph*)]
\item Show that the free group $F_2$ on $a,b$ represents $U^2$. Use Yoneda to prove that the operations $G\times G\to G$ natural in $G$ are precisely $(g,h)\mapsto w(g,h)$ for words $w\in F_2$.
\item Under the integer parametrisation $\operatorname{Nat}(U,U)\cong\Z$, find the parameter of inversion. Describe composition in terms of the parameters.
\end{enumerate}
\end{exercise}

\begin{exercise}[$\star$]\label{ex:matrixpoly}
For $n\ge0$, let $M_n:\mathbf{CRing}\to\mathbf{Set}$ send a ring to its $n\times n$ matrices and apply ring maps entrywise. Prove that $M_n$ is represented by $\Z[X_{ij}:1\le i,j\le n]$. Deduce that operations $M_n(R)\to R$ natural in $R$ are exactly integer polynomials in the entries, and give two familiar examples.
\end{exercise}

\begin{exercise}[$\star\star$]\label{ex:leftadj}
\begin{enumerate}[label=(\alph*)]
\item Prove that a left adjoint preserves coproducts, by applying Theorem~\ref{thm:rapl} in opposite categories.
\item Apply this to obtain $R^{(S\sqcup T)}\cong R^{(S)}\oplus R^{(T)}$ and $A\times(B\sqcup C)\cong(A\times B)\sqcup(A\times C)$ for sets. Describe the maps in both isomorphisms.
\end{enumerate}
\end{exercise}

\begin{exercise}[$\star\star$]\label{ex:f2plane}
Let $V$ be a two-dimensional vector space over $\mathbb F_2$. Define $\eta_V(v)(w)$ to be $1$ when $v,w$ are linearly independent, and $0$ otherwise.
\begin{enumerate}[label=(\alph*)]
\item Prove that each $\eta_V(v)$ is a linear form and that $\eta_V:V\to V^*$ is a linear isomorphism. Check \eqref{eq:dualnat} for every isomorphism between such planes.
\item Find the exact step in the proof of Theorem~\ref{thm:notnatural} that does not apply over $\mathbb F_2$.
\end{enumerate}
\end{exercise}

\begin{exercise}[$\star\star$]\label{ex:v10new-ch01-adjunction-order}
Fix an integer $d\ge1$ and regard $\Z$ as a category under its usual order. Define $F(n)=\lceil n/d\rceil$ and $G(k)=dk$.
\begin{enumerate}[label=(\alph*)]
\item Prove $F\dashv G$, including negative integers. Is $G\dashv F$ when $d>1$?
\item Compute the product and coproduct of two integers in this category. Check directly the product preservation of $G$ and the coproduct preservation of $F$.
\item Compare this example with $\gcd$ and $\operatorname{lcm}$ in the divisibility order. Why is the order relation part of the categorical data?
\end{enumerate}
\end{exercise}

\begin{exercise}[$\star\star$]\label{ex:v10new-ch01-abelian-operations}
Write abelian groups additively and let $U:\mathbf{Ab}\to\mathbf{Set}$ forget the addition.
\begin{enumerate}[label=(\alph*)]
\item Prove that every operation $U(A)\to U(A)$ natural in $A$ is $x\mapsto nx$ for a unique $n\in\Z$. Which are natural automorphisms of the identity functor on $\mathbf{Ab}$?
\item Classify the natural operations $U(A)\times U(A)\to U(A)$, and decide which of them have $0$ as a two-sided identity.
\end{enumerate}
\end{exercise}

\begin{exercise}[$\star\star$]\label{ex:v10new-ch01-subsets}
For a function $f:X\to Y$, let $\mathcal P(X),\mathcal P(Y)$ be the power sets ordered by inclusion. For $S\subseteq X$ and $T\subseteq Y$ define
\[
f_!(S)=f(S),\qquad f^{-1}(T)=\{x:f(x)\in T\},\qquad
f_*(S)=\{y:f^{-1}(\{y\})\subseteq S\}.
\]
Prove the two adjunctions $f_!\dashv f^{-1}\dashv f_*$. Use them to predict which maps preserve intersections or unions, then check those predictions directly, including empty intersections and empty unions. Explain why a point with empty fibre belongs to $f_*(\varnothing)$.
\end{exercise}

\begin{exercise}[$\star\star$]\label{ex:v10new-ch01-pullback}
Given functions $f:A\to C$ and $g:B\to C$, define
$P=\{(a,b)\in A\times B:f(a)=g(b)\}$, with its two projections.
\begin{enumerate}[label=(\alph*)]
\item Show that whenever $u:X\to A$, $v:X\to B$ satisfy $fu=gv$, there is exactly one $h:X\to P$ with the given components. This property defines a \emph{pullback}.
\item When $A,B,C$ are groups and $f,g$ homomorphisms, prove that $P$ is a subgroup and has the same property for homomorphisms.
\item For $A=B=\Z$, $C=\Z/2$, and both maps reduction modulo $2$, give a basis of the abelian group $P$. Decide whether $P$ is the cartesian product together with its usual projections.
\end{enumerate}
\end{exercise}

\section{Hints}

\begin{hint}{ex:zq}
For maps out of $\Q$, write $m/n$ using $m$ and the inverse of $n$. In sets, use two functions to $\{0,1\}$ that differ at one missing rational.
\end{hint}

\begin{hint}{ex:nofieldproduct}
A unital homomorphism between fields has zero kernel. What does this imply about their characteristics?
\end{hint}

\begin{hint}{ex:words}
A homomorphism $F_2\to G$ is specified by an arbitrary pair of images of $a,b$. For composition, evaluate a power operation on another power.
\end{hint}

\begin{hint}{ex:matrixpoly}
Use the universal property of a polynomial ring to encode a matrix by a homomorphism. The representing element supplied by Yoneda is itself an integer polynomial.
\end{hint}

\begin{hint}{ex:leftadj}
Reverse the arrows in the adjunction bijection: its inverse gives $G^{\mathrm{op}}\dashv F^{\mathrm{op}}$. Remember that the injections, as well as the objects, belong to a coproduct.
\end{hint}

\begin{hint}{ex:f2plane}
List the four vectors $0,u,v,u+v$. For nonzero $v$, the proposed form vanishes precisely on $\{0,v\}$.
\end{hint}

\begin{hint}{ex:v10new-ch01-adjunction-order}
For an integer $k$, $\lceil n/d\rceil\le k$ exactly when $n\le dk$. Test the reversed adjunction at $k=1,n=1$ if $d>1$.
\end{hint}

\begin{hint}{ex:v10new-ch01-abelian-operations}
The groups $\Z$ and $\Z^2$ represent one element and an ordered pair of elements. Evaluate an operation at $1$ or at the standard basis of $\Z^2$.
\end{hint}

\begin{hint}{ex:v10new-ch01-subsets}
Unpack $f(S)\subseteq T$ and $f^{-1}(T)\subseteq S$ one element at a time. A condition saying that all members of an empty set satisfy something is true.
\end{hint}

\begin{hint}{ex:v10new-ch01-pullback}
The only possible map is $h(x)=(u(x),v(x))$. In the last part, write a same-parity pair as $a(1,1)+b(0,2)$.
\end{hint}

\section{Solutions}

\begin{solution}{ex:zq}
\idea The category of rings restricts what a function out of $\Q$ can do.
\writeup The inclusion $i$ is injective, so $ig=ih$ implies $g=h$ pointwise: it is monic. Suppose ring maps $\varphi,\psi:\Q\to R$ agree on $\Z$. For every nonzero integer $n$, $\varphi(1/n)$ is the unique inverse of $\varphi(n)$, since $\varphi(n)\varphi(1/n)=1$ and the reversed product also equals $1$. Therefore
\[
\varphi(m/n)=\varphi(m)\varphi(n)^{-1}
 =\psi(m)\psi(n)^{-1}=\psi(m/n).
\]
So $\varphi=\psi$ and $i$ is epic. It is not invertible: a ring map $\Q\to\Z$ would make $2$ invertible in $\Z$. In sets, let $g,h:\Q\to\{0,1\}$ be zero everywhere except that $h(1/2)=1$. They agree after precomposition with $i$ but are unequal. Thus $i$ is not epic there. The same underlying inclusion has different cancellation behaviour because different outgoing maps are allowed.
\end{solution}

\begin{solution}{ex:nofieldproduct}
\idea Every allowed map between fields forces agreement of characteristics.
\writeup The kernel of a unital homomorphism $K\to L$ is an ideal of the field $K$ and does not contain $1$, hence is zero. Injectivity says that $n\cdot1_K=0$ if and only if $n\cdot1_L=0$. The two fields must have the same characteristic. A product would require maps $P\to\mathbb F_2$ and $P\to\mathbb F_3$, forcing both characteristics on $P$. A coproduct would require maps $\mathbb F_2\to C$ and $\mathbb F_3\to C$, with the same contradiction for $C$. The structural maps already fail to exist, before their universal properties can be tested.
\end{solution}

\begin{solution}{ex:words}
\idea A free group converts an arbitrary pair into the images of two fixed generators.
\writeup (a) The map $\operatorname{Hom}(F_2,G)\to U(G)^2$, $f\mapsto(f(a),f(b))$, is bijective by the free-group property. For $u:G\to H$, postcomposing $f$ gives $(u(f(a)),u(f(b)))$, so the bijection is natural. Hence $U^2\cong h^{F_2}$. Yoneda identifies $\operatorname{Nat}(U^2,U)$ with $U(F_2)=F_2$. A word $w$ yields $(g,h)\mapsto f(w)=w(g,h)$, where $f(a)=g$, $f(b)=h$. Distinct words give distinct operations, because evaluation at $(a,b)$ in $F_2$ distinguishes them. The word $ab$ gives multiplication and $aba^{-1}b^{-1}$ gives the commutator.

(b) Inversion has parameter $-1$. The composite $g\mapsto(g^m)^n=g^{mn}$ has parameter $mn$, including zero and negative parameters. Thus composition makes these natural operations the multiplicative monoid $(\Z,\times)$; it does not turn their parameters into the additive group of integers.
\end{solution}

\begin{solution}{ex:matrixpoly}
\idea A matrix is a choice of one ring element for each formal variable.
\writeup Put $P=\Z[X_{ij}:1\le i,j\le n]$. The assignment $f\mapsto(f(X_{ij}))$ is a bijection $\operatorname{Hom}(P,R)\to M_n(R)$. Applying a ring map entrywise is exactly postcomposition, so it is natural. Yoneda identifies natural operations to the underlying set of $R$ with elements of $P$. The element $Q\in P$ acts by evaluating $Q$ on the matrix entries. Trace and determinant are examples, as are coefficients of the characteristic polynomial. For $n=0$, $P=\Z$, there is a unique empty matrix, and the operations are integer constants; the trace and determinant are respectively $0$ and $1$.

This also explains the transfer of polynomial identities from complex matrices to arbitrary commutative rings. Every entry in an identity such as Cayley--Hamilton is an integer polynomial in the entries. An integer polynomial vanishing at all complex tuples is the zero polynomial: regard it as a polynomial in the last variable and use the fact that a nonzero one-variable polynomial has finitely many roots, then induct on the number of variables. Thus the identity holds after substitution in any commutative ring.
\end{solution}

\begin{solution}{ex:leftadj}
\idea Opposite categories turn outgoing universal maps into incoming ones.
\writeup (a) The inverse adjunction bijections can be written
\[
\operatorname{Hom}_{\mathcal C^{\mathrm{op}}}(G(B),A)
 \cong\operatorname{Hom}_{\mathcal D^{\mathrm{op}}}(B,F(A)).
\]
They are natural, so $G^{\mathrm{op}}\dashv F^{\mathrm{op}}$. The functor $F^{\mathrm{op}}$ is a right adjoint and preserves products by Theorem~\ref{thm:rapl}. A product in an opposite category is a coproduct in the original one, giving the assertion with its induced injections.

(b) The free-module functor is left adjoint to forgetting. Since $S\sqcup T$ is the coproduct of sets, its free module is the coproduct of $R^{(S)}$ and $R^{(T)}$. Explicitly a finitely supported function on the disjoint union restricts to one on each summand. Conversely the two functions combine on their disjoint supports. For sets, currying gives $A\times-\dashv(-)^A$, after swapping the factors. The resulting coproduct isomorphism sends $(a,b)$ to the $B$ summand and $(a,c)$ to the $C$ summand, preserving the tags distinguishing the two summands.
\end{solution}

\begin{solution}{ex:f2plane}
\idea Over $\mathbb F_2$, the kernel specifies a nonzero form on a plane completely.
\writeup (a) For nonzero $v$, the function $\eta_V(v)$ is zero on $\{0,v\}$ and one on the other two vectors. Choose $w$ independent of $v$. In the basis $(v,w)$, this is the second coordinate function, so it is linear. Also $\eta_V(0)=0$. Distinct nonzero vectors span distinct lines, hence give forms with distinct kernels. The map $\eta_V$ therefore bijects the three nonzero vectors with the three nonzero forms.

To check linearity in $v$, take distinct nonzero $v_1,v_2$. The form $\eta(v_1)+\eta(v_2)$ is nonzero, since its value at $v_1$ is $1$. Its value at $v_1+v_2$ is $1+1=0$. It is therefore the unique nonzero form vanishing on that vector, namely $\eta(v_1+v_2)$. If one vector is zero the identity is immediate; if they are equal, both sides are zero. Additivity implies $\mathbb F_2$-linearity. An isomorphism $f:V\to W$ preserves linear independence, so for every $v,w$,
\[
\eta_W(f(v))(f(w))=\eta_V(v)(w).
\]
This is precisely $f^*\eta_Wf=\eta_V$.

(b) The scalar argument required a nonzero $\lambda$ with $\lambda^2\ne1$. The only nonzero scalar in $\mathbb F_2$ is $1$, so it gives no contradiction. The example concerns planes and isomorphisms only. It does not provide a natural self-duality for all dimensions or for all linear maps.
\end{solution}

\begin{solution}{ex:v10new-ch01-adjunction-order}
\idea An integer lies below a ceiling exactly when it lies below the original number.
\writeup (a) Both maps preserve order, since $d>0$. For integers $n,k$,
\[
F(n)\le k\iff n/d\le k\iff n\le dk=G(k).
\]
The first equivalence holds for negative $n$ too, by the definition of the least integer above $n/d$. Thus $F\dashv G$. If $d>1$, $G(1)\le1$ is false whereas $1\le F(1)=1$ is true, so $G\dashv F$ fails.

(b) Arrows follow the ordinary order, so a product is $\min(a,b)$ and a coproduct is $\max(a,b)$. Multiplication by the positive integer $d$ commutes with minimum. The increasing map $F$ commutes with maximum: if $a\le b$, both sides of $F(\max(a,b))=\max(F(a),F(b))$ equal $F(b)$. This verifies the two preservation claims without an abstract argument.

(c) With divisibility instead, common sources are common divisors and common targets are common multiples, giving gcd and lcm. An arrow means a different relation, so even the same underlying objects need not have the same universal constructions.
\end{solution}

\begin{solution}{ex:v10new-ch01-abelian-operations}
\idea Use the universal examples $\Z$ and $\Z^2$, then map their generators wherever desired.
\writeup (a) For a natural operation $\eta$, write $n=\eta_{\Z}(1)$. For $x\in A$, the homomorphism $f:\Z\to A$ with $f(1)=x$ gives $\eta_A(x)=f(n)=nx$. Conversely multiplication by any integer commutes with every homomorphism. It is itself a homomorphism, so these are the natural endomorphisms of the identity functor as well. To be an automorphism for every $A$, it must be an automorphism on $\Z$, forcing $n=1$ or $-1$. Both choices work on every abelian group.

(b) Let $e_1,e_2$ be the basis of $\Z^2$. The value of a natural binary operation there is $me_1+ne_2$ for unique integers $m,n$. A map $\Z^2\to A$ carrying $e_1,e_2$ to $x,y$ forces the value $mx+ny$. Every such formula is natural. Having $0$ as a right identity requires $mx=x$ for all groups and elements, hence $m=1$ by testing $\Z$; a left identity forces $n=1$. Thus addition is the only binary operation in this class with two-sided identity $0$.
\end{solution}

\begin{solution}{ex:v10new-ch01-subsets}
\idea Existential and universal conditions on fibres are the two adjoints to inverse image.
\writeup For $S\subseteq X,T\subseteq Y$,
\[
f(S)\subseteq T\iff S\subseteq f^{-1}(T),
\qquad
f^{-1}(T)\subseteq S\iff T\subseteq f_*(S).
\]
The first equivalence says that every $x\in S$ has $f(x)\in T$. The second says that for every $y\in T$, all points mapping to $y$ belong to $S$. These prove the two adjunctions.

Products in power-set orders are intersections and coproducts are unions. Right adjoints preserve intersections; left adjoints preserve unions. For arbitrary indexed families, direct checking gives
\[
f_!\Bigl(\bigcup_iS_i\Bigr)=\bigcup_i f_!(S_i),\qquad
f^{-1}\Bigl(\bigcup_iT_i\Bigr)=\bigcup_i f^{-1}(T_i),
\]
\[
f^{-1}\Bigl(\bigcap_iT_i\Bigr)=\bigcap_i f^{-1}(T_i),\qquad
f_*\Bigl(\bigcap_iS_i\Bigr)=\bigcap_i f_*(S_i).
\]
For example, a fibre is contained in every $S_i$ exactly when it is contained in their intersection. With no sets in the family, the union is empty and the intersection is the entire ambient set. The four identities still hold: $f_!(\varnothing)=\varnothing$, $f^{-1}(\varnothing)=\varnothing$, $f^{-1}(Y)=X$, $f_*(X)=Y$. Finally an empty fibre is contained in $\varnothing$, so every point outside $f(X)$ belongs to $f_*(\varnothing)$. This is required by the universal condition, rather than an exception to it.
\end{solution}

\begin{solution}{ex:v10new-ch01-pullback}
\idea A compatible pair of maps already gives the only possible map to the subset of pairs.
\writeup (a) Define $h(x)=(u(x),v(x))$. The identity $fu=gv$ ensures $h(x)\in P$. Its projections are $u,v$. Any map with those projections has that ordered pair as its value at every $x$, proving uniqueness.

(b) The identity pair belongs to $P$. If $(a,b),(a',b')\in P$, then $f(aa')=f(a)f(a')=g(b)g(b')=g(bb')$, so the product is in $P$; inverses satisfy the same equality. Thus $P$ is a subgroup of $A\times B$. Its projections are homomorphisms. If $u,v$ are homomorphisms, the formula for $h$ respects multiplication componentwise, so the same existence and uniqueness proof works in $\mathbf{Grp}$.

(c) Here $P$ is the group of same-parity pairs $(a,b)$. Each such pair is uniquely
\[
(a,b)=a(1,1)+\frac{b-a}{2}(0,2).
\]
Consequently $(1,1),(0,2)$ is a basis and $P\cong\Z^2$ as an abstract group. It is a proper subgroup of the cartesian product, since $(0,1)$ is absent. With its specified projections it is therefore not the product of $\Z$ with itself: the pair of maps taking $1$ to $0$ and $1$ would require a lift with value $(0,1)$. The distinction is between an abstract isomorphism of objects and an isomorphism preserving the structural maps.
\end{solution}

\section*{Sources and further reading}
\addcontentsline{toc}{section}{Sources and further reading}

S.~Eilenberg and S.~Mac Lane, ``Group extensions and homology'', \emph{Ann.\ of Math.} 43 (1942), 757--831, and ``General theory of natural equivalences'', \emph{Trans.\ Amer.\ Math.\ Soc.} 58 (1945), 231--294. The story of the
Michigan lecture and Steenrod's solenoid is told by Mac Lane himself in H.~Bass, H.~Cartan, P.~Freyd, A.~Heller and S.~Mac Lane, ``Samuel Eilenberg 1913--1998'', \emph{Biographical Memoirs of the National Academy of
Sciences} 79 (2001). The meeting at the Gare du Nord is recalled by Y.~Kinoshita in his note on Yoneda's death, posted to the categories mailing list on 23 April 1996. D.~M.~Kan, ``Adjoint functors'', \emph{Trans.\ Amer.\ Math.\ Soc.} 87 (1958),
294--329. Mac Lane's summary of the logic (categories in order to define functors, functors in order to define natural transformations) and his
slogan ``adjoint functors arise everywhere'' are both in \emph{Categories for the Working Mathematician}: the slogan in the preface, the summary in Section~I.4.

E.~Moggi, ``Notions of computation and monads'', \emph{Information and Computation} 93 (1991), 55--92; P.~Wadler, ``Monads for functional programming'', in \emph{Advanced Functional Programming}, Lecture Notes in Computer
Science 925 (Springer, 1995).

For more: S.~Mac Lane, \emph{Categories for the Working Mathematician} (Springer, 2nd ed.\ 1998), Chapters~I--IV; T.~Leinster, \emph{Basic Category Theory} (Cambridge University Press, 2014), short and very clear, freely available
from its author; E.~Riehl, \emph{Category Theory in Context} (Dover, 2016), with examples from every part of mathematics.
